\section{Exploration Strategies}
\label{sec:exploration}

Exhaustive evolution becomes infeasible at moderate depth, and a caller may need only part of
it. The engine has four independent ways to limit an exploration, and they can be combined:

\begin{itemize}
\item \textbf{Caps.} A total state count, a total event count, a count of states admitted per step
      (\textsc{MaxStatesPerStep}) and a count of successors per parent
      (\textsc{MaxSuccessorStatesPerParent}). Each is an atomic counter, so the count retained is
      exact at every worker count; which states fill a cap depends on arrival order
      (Section~\ref{sec:determinism}). A run that reaches a cap returns what it computed with a
      warning (Section~\ref{sec:capacity}).
\item \textbf{Sampling.} Per-rule transition rates and a per-state match cap keep a subset of
      transitions chosen by keyed pseudorandom draws, so the sample is reproducible for a seed at
      every worker count (Section~\ref{sec:sampling}).
\item \textbf{Quotient exploration} (Section~\ref{sec:quotient}) expands one state per isomorphism
      class and reconstructs the raw relations exactly.
\item \textbf{Steering.} A session (Section~\ref{sec:sessions}) reports the frontier after each
      step, and the next step can be restricted to part of it; the rest stays on the frontier.
\end{itemize}

Caps and steering compute a subset exactly; sampling computes a random subset; quotient exploration
computes the full result and expands less. The step budget bounds all four, and the reported
frontier records where each stopped.
\label{sec:pruning}

\subsection{Quotient Exploration}
\label{sec:quotient}

In full expansion every raw state is expanded, so a state reached along $k$ histories is matched $k$
times, although canonicalization has already shown the repeats find no new states. Quotient
exploration expands only one state per isomorphism class, at the shallowest depth the class was
reached. Every equivalent state is still created and its incoming event recorded.
Table~\ref{tab:t6-quotient} gives events produced, causal edges and wall time for both modes by
depth. Both modes give the same canonical state set. The causal and branchial relations are
defined over raw events, which quotient exploration does not all produce, so they are
reconstructed.

The number of raw events is recovered by a forward summation. Under exact identity the canonical
states graph has cycles, because a canonical state recurs at several depths, but its unrolling by
depth is acyclic, and the multiplicity of each (canonical state, depth) pair satisfies
\begin{align}
  \mathrm{mult}[s_0, 0] &= 1, \nonumber\\
  \mathrm{mult}[c, k{+}1] &\mathrel{+}= \mathrm{mult}[s, k]\cdot w(s \to c),
  \label{eq:mult}
\end{align}
summed over the transitions into $c$; the raw event count is $\sum \mathrm{mult}$. This takes time
polynomial in the number of (canonical state, depth) pairs.

The raw events and relations themselves are recovered by replaying each class's expansion against
every instance of the class. The expansion is recorded once, in the \emph{frame} of the first raw
state that reached the class: each match found on that state is stored as indices into the frame's
edges. Another raw state of the class is aligned onto the frame by the canonical edge ranks from
its own IR run; this alignment is an isomorphism, defined up to an automorphism, so it maps matches
to matches. The edge orbits from the same IR pass distinguish edge roles within a class
(Section~\ref{sec:ir}), so no further canonicalization is needed. An \emph{instance} of a class at
depth $k$ is a vector holding one producer event per frame slot. Applying a recorded match to an
instance creates one raw event; adds a causal edge from the producer of each consumed slot, in
descending producer order so that the online reduction of Section~\ref{sec:tr} stays exact; adds a
branchial pair with each earlier application at the same instance whose consumed slots overlap;
and creates a child instance at depth $k+1$, in which surviving slots keep their producer and
produced slots take the new event. Instances and recorded matches arrive concurrently, so each new
instance replays every match already recorded for its class, each new match is replayed against
every instance already present, and each (instance, match) pair is claimed exactly once. A replayed
event has two identities: the triple (source class, target class, rule), which is an isomorphism
invariant used for comparisons across runs and engines, and the run's event identity.

The reconstruction creates as many instances as the full system has raw states. On a two-rule set
whose canonical state count grows linearly with depth (three to eight states at depths two to
seven), the raw system has 146,599 states at depth six, and the reconstruction's cost follows that
count. There, quotient exploration visited 49 raw states against full capture's 146,599 and ran 65
times faster.

The reconstruction runs only when the causal relation, the branchial relation or the raw event set
is requested; a run that asks for canonical states alone skips it. Every worker takes part in it,
and its total cycle count grows with the thread count (Section~\ref{sec:limitations}).

\subsection{Sampled Exploration}
\label{sec:sampling}

The engine has two samplers. The \emph{transition rate} keeps each candidate transition
independently with a fixed probability, multiplied by a per-rule weight; the number of children a
state keeps is then a random variable with the thinned distribution of the full one. The
\emph{per-state match cap} keeps at most $k$ of each state's matches for each rule. Both choose by a
pseudorandom value computed from the transition's canonical identity and the seed, so the same seed
gives the same sample at any thread count (Section~\ref{sec:determinism}).

\paragraph{Choosing the rate.}
A fixed rate compounds with depth: a path of length $L$ survives with probability $q^L$. On the
workloads measured, rates of $1/8$ and $1/4$ end the evolution within the first step and $1/2$
leaves it growing, so the usable range is narrow and depends on the rule's branching factor. While
a rate or weight below one is in effect, a state whose own draws all fail keeps its
lowest-ranked own transition, so the evolution does not die out; a state with no matches of its
own has no such fallback.

\paragraph{The per-state cap.}
Samples from different seeds do not add up to the exhaustive system, because a transition can be
drawn only if its source state was created, which required an earlier draw to succeed. The
per-state cap keeps at least one transition at every state it reaches, so each seed covers a
connected skeleton, and different seeds keep different skeletons. The cap chooses $k$ of a state's
$M$ matches by rank once the state's matching is complete, because choosing $k$ of $M$ needs all
$M$; choosing by arrival order would make the result depend on the schedule. A match forwarded
from an ancestor can arrive after the state's matching is complete, so a run under the cap turns
match forwarding off and each state finds all of its matches itself.

\paragraph{Per-state and per-transition draws.}
Under full capture each raw state is created by exactly one event, so a draw per new state and a
draw per transition are the same draw. Under quotient exploration a canonical state has many
incoming transitions, and the two differ.

\paragraph{No barrier.}
A decision to keep a match is final once the match is rewritten, so a sample must be complete
before its members are rewritten. A sample drawn from a whole evolution step would be complete only
when every worker finished the step, which is a global barrier. The engine therefore samples only
over populations that complete locally: a state's matches for one rule are complete when that
state's matching tasks have all finished, which the engine detects without looking at any other
state.

\subsection{Sessions}
\label{sec:sessions}

Evolution is also available as a session: it is opened on a rule set and an initial state, advanced
by a number of steps, queried for any derived structure, and closed. The frontier is kept between
steps, so each step continues the existing evolution instead of starting again.

The frontier is stored separately from the canonical identity map. The map records which states
have been seen; the frontier records the states and pending rewrites that the step budget stopped,
together with the depth each waits at.

Every session reply reports the frontier, and a step may name part of it. The named entries are
advanced and the others stay on the frontier. A step that names no entries advances the whole
evolution to the new total depth, and gives the same result as a single run to that depth. A step
that names entries advances each of them by the requested number of steps from the depth it waits
at, so entries at different depths each advance by the same amount. Naming a state that is not on
the frontier is an error. A continuable run stopped by a state or event limit keeps the work the
limit cut short on the frontier (Section~\ref{sec:capacity}). Both devices implement this interface.

